\section{The Results \label{sec:theresults}}

We will now compare our results with BKK's in order to address the question of whether the empirical regularities identified in 1993, specifically those regarding international comovements, are still valid. To make comparison easier, our results are presented beside BKK's, which appear in parenthesis.

\subsection{Within-country Comovements and Volatility}

We report the values for standard deviations and within-country correlations in Table \ref{tab:stdv} and Table \ref{tab:corrwithin}, respectively. Table \ref{tab:stdv} gives information about volatility by displaying the ratio of the standard deviation of each variable to the standard deviation of output (except for net exports, whose standard deviation is presented in absolute terms, following BKK). Table \ref{tab:corrwithin} shows comovements of economic variables inside each country. Its second column shows the autocorrelation of GDP with its one-period-lagged value and the following columns report the contemporaneous correlation between GDP and the respective variable.


%Standrd Deviation
% \begin{table}[htb]
% \centering
% \begin{tabular}{Hlcccccc}
%   \hline \hline
%    & & \multicolumn{2}{c}{Std. Dev. (\%)} & & \multicolumn{3}{c}{Ratio of Std. Dev. to that of y}\\ \cline{3-4} \cline{6-8}
%  & Country & y & nx & & c & x & g \\
%   \hline
%   1 & Australia     & 1.32 (1.45) & 2.68 (1.23) & & .83 (.66) & 2.98 (2.78)& 1.27 (1.28)\\
%   2 & Austria       & 1.24 (1.28) & 2.55 (1.15) & & .94 (1.14)& 2.65 (2.92)& .68 (.36)\\
%   3 & Canada       & 1.36 (1.50)& 2.79 (.78)   & & .81 (.85)& 2.93 (2.80)& 1.04 (.77) \\
%   4 & European Union   & 1.07 (1.01) & 2.84 (.50)   & & .72 (.83)& 2.44 (2.09)& .53 (.47)\\
%   5 & France       & 1.12 (.90) & 2.44 (.82)   & & .79 (.99)& 2.30 (2.96)& .58 (.71)\\
%   6 & Germany      & 1.58 (1.51) & 1.02 (.79)  & & .77 (.90)& 2.57 (2.93)& .86 (.81)\\
%   7 & Italy       & 1.46 (1.69) & 1.13 (1.33) & & .87 (.78) & 2.55 (1.95)& .58 (.42)\\
%   8 & Japan       & 1.56 (1.35)& 1.79 (.93)   & & .79 (1.09) & 2.22 (2.41)& .74 (.79)\\
%   9 & Switzerland     & 1.66 (1.92) & 1.63 (1.32) & & .64 (.74) & 2.36 (2.30)& .92 (.53)\\
%   10 & United Kingdom   & 1.48 (1.61) & 1.50 (1.19) & & 1.08 (1.15) & 2.58 (2.29)& .84 (.69) \\
%   11 & United States   & 1.63 (1.92) & 1.63 (.52)  & & .77 (.75) & 2.48 (3.27)& .85 (.75) \\
%    \hline \hline
% \end{tabular}
% \end{table}
\input{../output/tables/table_4_standard_deviations.tex}

% % latex table generated in R 3.1.2 by xtable 1.7-4 package
% % Sat Aug 29 19:22:56 2015
% \begin{table}[htb]
% \centering
% \begin{tabular}{HlcHcccccc}
%   \hline
%  & Country & Autocor & y & c & x & g & nx & n & z \\
%   \hline
% 1 & AUS & .70 & 1 & .44 & .66 & .00 & -.11 & .56 & .84 \\
%   2 & AUT & .72 & 1 & .54 & .68 & -.05 & .01 & .29 & .91 \\
%   3 & CAN & .83 & 1 & .78 & .69 & -.21 & -.04 & .76 & .89 \\
%   4 & CHE & .83 & 1 & .70 & .78 & .34 & -.03 &  &  \\
%   5 & DEU & .78 & 1 & .63 & .84 & -.06 & .03 & .58 & .91 \\
%   6 & EU15 & .86 & 1 & .80 & .91 & -.20 & -.16 &  &  \\
%   7 & FRA & .55 & 1 & .68 & .83 & -.03 & -.13 & .58 & .93 \\
%   8 & GBR & .80 & 1 & .77 & .74 & -.08 & -.14 & .62 & .92 \\
%   9 & ITA & .82 & 1 & .69 & .79 & -.04 & -.17 & .33 & .95 \\
%   10 & JPN & .78 & 1 & .75 & .86 & .07 & -.06 & .46 & .98 \\
%   11 & USA & .85 & 1 & .77 & .92 & -.13 & -.11 & .77 & .92 \\
%    \hline
% \end{tabular}
% \end{table}
\input{../output/tables/table_5_within_country_correlations.tex}


Our results are in line with what have come to be known as the stylized facts of the business cycle.
Consumption is procyclical and, except for UK, less volatile than output. Differently from BKK's result, in which three countries had consumption more volatile than output, consumption volatility relative to output seems to have decreased. From the eleven countries, eight have experienced reductions in this measure. From the three that have seen its consumption volatility rise, only Australia experienced an increase greater than 0.1.

Investment is also procyclical but more volatile than output. In Japan, the country with the lowest volatility, investment has been 2.22 times more volatile than output, while in Australia, where it is most volatile, investment has been almost three times more volatile than output. Different from consumption, there has not been a clear tendency in the behavior of investment volatility from BKK years to nowadays. The sharp decrease US has experienced in investment volatility is worth noting, however. Standard deviation of investment relative to that of output in the US has decreased from 3.27 to 2.48. This combined with the Italian increase has led to a reduction of the range of the investment volatility distribution from 1.32 in BKK to 0.72.

Autocorrelation of output has not changed in any noticeable way. Six countries have witnessed an increase and five a decrease in this measurement, most of which of modest magnitude. Output volatility has decreased for six countries and increased for five. Remarkable, however, is that the distribution has become more uniform. Countries that had higher volatilities of output have in general experienced a decrease (United States, Switzerland and Italy), while those with lower volatilities have seen their volatilities increase (France, European aggregate and Japan). The tails of the distribution are now closer to the mean.

Government spending does not display a clear tendency regarding its correlation with output. In our sample most of the countries have a negative correlation between output and government spending but with small magnitude. Moreover, the largest correlation in magnitude is positive (Switzerland, 0.34). Together, these two facts -- majority of small-magnitude negative correlations and the highest correlation in magnitude being positive -- makes a strong case for the acyclicality of government spending. In BKK's data set, correlations between output and government spending are most positive but never greater than 0.27 in magnitude, which also indicates that government spending has being acyclical since then.

The volatility of employment and its correlation with output have changed considerably for some countries but the overall picture have stayed constant. The same happened to Solow residuals: some countries experienced increases in its volatility, while others have seen it decline. The overall distribution of the cross-country correlations of both employment and Solow residuals, however, does not seem to have changed.

\subsection{International Comovements of the USA}
Having seen that stylized facts within economies have barely changed, we now turn to the comovements among countries. We use -- as BKK did -- the US economy as a benchmark and measure how its economic variables correlate with the respective variables of the other ten countries in the sample. Table \ref{tab:corusa} presents the contemporaneous cross-country correlations between USA and the other ten countries.

In accord with BKK findings, outputs are positively correlated. The magnitude of the correlations, however, seems to have decreased. Of ten output correlations, nine have become smaller. Only UK's output experienced a slight increase in its correlation with American output. Smaller cross-country correlations of output point in the direction of an approximation between theory and data.

Consumption behaves similarly as in BKK. In the longer time frame of our data, there is a positive correlation between the US and all countries but Italy. In BKK's data set, US consumption correlation was positive with all countries except for Australia. Although not as clear as in the case of output, there seems to be tendency for the magnitude of consumption correlation to decrease. Different from the decrease in cross-country correlation of output, the decrease in consumption correlation across countries accentuates the contrast between theory and data, since the model predicts very high correlation of consumptions (0.88).

Comovements of investment and government spending between US economy and the other countries have not changed their behavior in any remarkable way. Nonetheless, we note that, although US investment still correlates positively with the other sample countries', the range of the distribution have decreased from 0.57 to 0.39. Cross-country correlations of government spending are still very low meaning that government spending is virtually uncorrelated between countries.

From our calculations, net exports appear to be uncorrelated between US and the other countries. If anything, a light tendency to a negative correlation may be identified. BKK have not reported cross-country correlations of net exports.

US employment correlates positively with other countries' in our data as well as in BKK's. The cross-country correlation of Solow residuals is still positive between the US and all other countries but its magnitude seems to have decreased.

Updating BKK's calculations with 23 years of recent data does not change the conclusions they have arrived at in 1993. Cross-correlations of output and consumption might have decreased on average. Yet, this fact does not do much to disprove the central discrepancy BKK had found between theory and data. What characterizes the quantity anomaly is not the output or consumption correlation itself but precisely the relative magnitude of both. Since consumption and output correlations have decreased together, the relation between them has stayed virtually the same and so too the quantity anomaly.

\renewcommand{\baselinestretch}{1.3}
% \begin{table}[htb]
% \centering
% \begin{tabular}{lccccc}
%   \hline \hline
%  & \multicolumn{5}{c}{Correlation with Same U.S. Variable}\\\cline{2-6}
% Country & y & c & x & g & nx \\
%   \hline \hline
%   AUS & .27 (.51) & .02 (-.19) & .36 (.16) & .07 (.23) & .04 \\
%   AUT & .28 (.38) & .17 (.23) & .33 (.46) & .09 (.29) & -.74 \\
%   CAN & .75 (.76) & .56 (.49) & .35 (-.01) & .33 (-.1) & -.18 \\
%   CHE & .30 (.42) & .21 (.40) & .38 (.38) & -.03 (.01) & -.46 \\
%   DEU & .37 (.69) & .30 (.49) & .42 (.55) & .05 (.28) & -.69 \\
%   EU15 & .53 (.66) & .44 (.51) & .51 (.53) & .23 (.18) & -.37 \\
%   FRA & .36 (.41) & .29 (.39) & .35 (.22) & -.03 (-.20) & -.11 \\
%   GBR & .59 (.55) & .45 (.42) & .44 (.40) & .21 (-.04) & .78 \\
%   ITA & .32 (.41)& -.03 (.02) & .12 (.31) & .14 (.09) & -.03 \\
%   JPN & .32 (.60) & .22 (.44) & .35 (.56) & .03 (.11) & -.49 \\
%    \hline \hline
%      \end{tabular}
% \caption{International Comovements in BKK's Countries \label{tab:corusa}}
% \end{table}

\input{../output/tables/table_6_us_correlations.tex}



\subsection{International Comovements of all Countries}

While analyzing international comovements, BKK have only used correlations involving the US economy. By doing so, they have ignored other 45 cross-country correlations for each variable resulting from the sample with eleven countries. As \cite{Ambler:2004aa} put, the US economy is not representative among industrialized economies in terms neither of size nor of openness. For that reason, cross-country correlations of US economy might not provide a good picture of what happens among the other countries.

We inquire next how well the correlations involving the US represent the distribution of all correlations. We are interested here not only in how United States' output correlates with, say, Germany's and Japan's but also in how these two correlate with each other. To address this issue, we plotted in Figure \ref{fig:cor} all the 55 correlations of output and consumption in ascending order. Output and consumption were chosen for being the two most relevant variables for the quantity anomaly. Correlations involving the US are highlighted, so that the reader can have an idea of the implications of choosing the US as the benchmark for determining the properties of international comovements.

Output correlations are all positive. Those involving US are on average smaller than the others. Consumptions correlations are also mainly positive with two exceptions (ITA-USA and AUS-DEU). For consumption, however, US correlations are more evenly spread along the distribution.

\begin{figure}[htb]
\centering
\subfloat[Output Correlations]{
\includegraphics[width=0.45\textwidth]{../output/figures/gdpcor.png}
\label{fig:gdpcor}}
\qquad
\subfloat[Consumption Correlations]{
\includegraphics[width=0.45\textwidth]{../output/figures/concor.png}
\label{fig:concor}}
\caption{Output and Consumption Correlations among all Countries}
\label{fig:cor}
\end{figure}

We now approach the problem of determining the characteristics of the international comovements from a perspective other than BKK's. Instead of reporting only the cross-country correlations of the US, we calculate the mean of the correlations between all countries of our sample. We also calculate the mean of the correlations involving the US to assess in how far BKK's approach deviates from ours. \cite{Ambler:2004aa}, hereafter ACZ, adopt a similar approach\footnote{ACZ also reassess BKK's results using a greater sample of countries and slightly different data. They use a sample of twenty industrialized countries with data from 1960:1 to 2000:4. Different from BKK, they  include capital when calculating the Solow residuals. They also report the standard errors and test for the statistical significance of their results.} and we report their results in square brackets. Table \ref{tab:avgcor} shows in its first column the average of all correlations, in the second column the average correlations involving US and in the remaining columns various quantiles summarizing the distribution of all correlations.

% \begin{table}[ht]
% \centering
% \begin{tabular}{r>{\bfseries}ccrrrr>{\bfseries}crrrr}
%   \hline
%  & Mean & USA Mean & 0\% & 10\% & 25\% & 40\% & 50\% & 60\% & 75\% & 90\% & 100\% \\
%   \hline
% y & 0.45 & 0.41 (0.54)& 0.06 & 0.23 & 0.31 & 0.39 & 0.44 & 0.48 & 0.60 & 0.71 & 0.86 \\
%   c & 0.26 & 0.26 (0.32) & -0.22 & 0.05 & 0.12 & 0.17 & 0.25 & 0.29 & 0.39 & 0.56 & 0.70 \\
%   x & 0.38 & 0.36 (0.36) & 0.01 & 0.15 & 0.24 & 0.32 & 0.37 & 0.42 & 0.51 & 0.61 & 0.83 \\
%   g & 0.10 & 0.11 (0.35) & -0.28 & -0.07 & -0.03 & 0.05 & 0.07 & 0.11 & 0.15 & 0.33 & 0.78 \\
%    \hline
% \end{tabular}
% \end{table}
\input{../output/tables/table_7_average_cross_country_correlations.tex}


We first note that the mean and the median of all correlations are very similar, so they can be used interchangeably to characterize the distribution. The quantiles of the distribution also indicate that the quantity anomaly is unlikely to have disappeared. 50\% of the output correlations are greater than 0.44, while none of the consumption correlations is as high as 0.88, the expected value in theoretical economies. The largest consumption correlation is 0.70 and 90\% of them lie below 0.56.


It turns out that averaging only USA correlations delivers surprisingly similar figures to the ones resulting from averaging all correlations. If net exports are not considered,\footnote{US trade balance has a \textit{sui generis} behavior due to the status of the dollar as the international reserve currency. The rate of return on US foreign investments is two to three times higher than the interest rate the US has to pay for liabilities of other countries. This allows the US to constantly run trade deficits without becoming more indebted. For that reason, US comovements of net exports might not be a good proxy for the whole distribution of cross-country correlation of net exports. For more on what have become to be known as the exorbitant privilege, see \cite{Eichengreen:2011aa}} the largest difference between overall and USA mean is in civilian employment, which amounts to 0.07. Except for investment and employment, the correlations of variables have decreased when compared to BKK's results. USA mean of output correlations have dropped by 0.13 and consumption by 0.06. However, these decreases are not as striking as those reported by ACZ. Output correlation drops from 0.54 in BKK to 0.22 in ACZ. In this respect, ACZ's results brings the data closer to the model, which predicts low and, depending on the version of the model, even negative cross-correlation of output. Consumption, on the other hand, decreases by 0.18, a result which accentuates the difference between model and data, since cross-country correlation of consumption should be very high according to the theory.

Our results represent a middle ground between BKK and ACZ. The greater sample used by ACZ seems to be key for the large decrease in all the cross-country correlations. However not as striking as in ACZ, the larger time frame we have adopted also had an impact in diminishing the correlations among countries.

This overall decrease in cross-country correlations have little impact on the quantity anomaly.On the one hand, the overall decrease in cross-country correlations attenuates the discrepancy between theory and data, because it means that output correlations will become closer to the levels predicted by the model. On the other hand, the decrease in cross-country correlation of consumption accentuates this discrepancy, since consumption correlation is very high in the theoretical economy.

In fact, the quantity anomaly becomes even more salient when each pair of cross-country correlation is considered individually. In Figure \ref{fig:gdpconplot} we plot the cross-country correlation of output and consumption for each of the 55 possible pairs of countries. The line has a slope of 45 degrees and represent all combinations of values for which output correlation equal consumption correlation. Of all pairs, only two lay below the line and thus violate the empirical regularity that outputs correlate more strongly than consumptions (AUS-FRA and CHE-GBR in red). For all the others, correlation of outputs is greater than correlation of consumption.

\begin{figure}[htb]
\centering
\includegraphics[scale=.15]{../output/figures/gdpconplot.png}
\caption{Cross-country Correlation of Output and Consumption for each Pair of Countries \label{fig:gdpconplot}}
\end{figure}

In sum, the quantity anomaly persists: the difference between output and consumption correlation is roughly 0.20 both in BKK's calculations and ours. Although ACZ arrives at a milder difference of only 0.08, they reject the null hypothesis that the cross-correlation of output is less than or equal to the average cross-correlation of consumption ($H_0: \rho_{y,y^*} \leq \rho_{c,c^*}$) at 5\% significance level. What's more, when the pairs of countries are taken into consideration individually, the quantity anomaly becomes even more evident: out of 55 pairs, 53 have correlation of output that are greater than that of consumption.


\subsection{Further Research on International Business Cycles}

We indicate next two papers that claim to have corrected the quantity anomaly, at least partially, and point out different directions that research in international business cycle have taken in the last years.

\cite{Kehoe:2002aa} corrects the problem of negative cross-country correlation of output. They build a model in which countries can default on their loans: the single mechanism to enforce loans is the threat of excluding the defaulting country from future trade. This friction in the credit market generates a positive cross-country correlation of output (0.25), which is still smaller than that of consumption (0.29), however.

\cite{Faia:2007aa} disputes the idea that the central driver of international comovements is trade and makes the case for financial institutions as a good predictor of the size of the comovements between countries. She shows that there is a negative relation between pairwise output correlation and a measure of financial distance and constructs a model with financial frictions that mimics the positive comovements of output, investment and employment.

Other works have taken different approaches in analyzing the international business cycles. \cite{Kose:2003aa} have used a Bayesian method on yearly data to divide the international business cycle into three factors (world, region and country-specific factors) and have come to the conclusion that the common world factor plays and important role in explaining the business cycles for a 60-country sample. \cite{Stock:2003aa} apply a structural vector autoregression (VAR) model\footnote{\cite{Webb:1999aa} provides a simple and enlightening introduction to VAR models.} to explain the reduction of volatility among G7 countries in the last two decades and conclude that the less volatile economies result from the decrease of the common international shocks. \cite{Fund:2010aa} apply also a structural VAR to highlight the relevance of US financial shocks to a main driver of international business cycles.

There is a variety of approaches that can be adopted while studying international business cycles. In the recent years, it seems that more advanced econometric techniques, such as structural VAR models, have been more often used than the simple analysis of comovements among countries, which we have adopted in this paper.
